\documentclass[11pt]{article}

\usepackage[margin=1in]{geometry}
\usepackage[T1]{fontenc}
\usepackage[utf8]{inputenc}
\usepackage{lmodern}
\usepackage{microtype}
\usepackage{amsmath,amssymb}
\usepackage{booktabs}
\usepackage{enumitem}
\usepackage[hidelinks]{hyperref}
\usepackage{xcolor}

\setlist{itemsep=2pt, topsep=4pt}

\title{\textbf{Diffusal: Error Dynamics of Ternary Quantization\\in Masked Diffusion Language Models}\\[6pt]
\large A Research Program Proposal}
\author{Victor Panisa\\\texttt{victor.panisa@gmail.com}}
\date{July 2026}

\begin{document}
\maketitle

\begin{abstract}
Diffusion large language models (dLLMs) refine a token canvas iteratively, applying the same denoiser repeatedly to (partially) its own output. This makes them a distinctive target for extreme weight quantization: weight error can alter both a refinement state and the token or confidence decisions that determine later states. We propose \emph{Diffusal}, a measurement-then-build program whose end goal is a natively ternary ($\{-1,0,+1\}$, 1.58-bit) masked text diffusion model, obtained by autoregressive-to-diffusion (A2D) conversion of a small open-weight model followed by ternary quantization-aware training (QAT). The central scientific question is whether ternary quantization causes more, equal, or less quality loss in iterative refinement than in a matched autoregressive (AR) model, and which sampler pathways explain the result. A crossover in state stability across mask fractions is a falsifiable working hypothesis, not an assumed property of the model. Unlike continuous image diffusion, where step-dependent error accumulation is established, masked text diffusion has a mechanism with no continuous analog: a deterministic \textbf{token-commit decision} can suppress a sub-margin logit perturbation at one position, while a margin crossing produces a discrete, often irreversible token flip. Confidence and remasking decisions remain possible pathways even when token identity is unchanged. We propose a three-experiment program with thresholds committed before data collection, a paired perturbation-stability protocol that separates conditional distortion from state recovery under each transition map, and a matched autoregressive control at every precision. A dLLM-only ternary failure cannot identify diffusion as the cause; only excess degradation relative to that control is evidence about diffusion.
\end{abstract}

\section{Introduction}

Masked and discrete diffusion language models such as MDLM~\cite{mdlm} and LLaDA~\cite{llada} generate text by iteratively denoising a token canvas rather than appending tokens left to right. In parallel, BitNet b1.58~\cite{bitnet} argues that ternary-weight LLMs can be competitive with full-precision models when trained natively at that precision. Neither result implies the other's applicability: whether iterative refinement tolerates, absorbs, or amplifies ternary quantization noise is an open question, and it is the question this program is built around.

The engineering motivation is concrete. A diffusion large language model (dLLM) applies its denoiser $T$ times per generation; an autoregressive (AR) model applies its network once per token but never revisits emitted tokens. If refinement dynamics are \emph{contractive}, a quantized dLLM may be unusually robust---each step corrects the previous step's error. If they are \emph{expansive}, quantization error compounds in a way AR generation structurally cannot, and low-bit dLLMs are more fragile than their weight count suggests. The realistic answer, we argue, is a \textbf{crossover}: contraction at high mask fractions and expansion near convergence. The primary deliverable is the location of that crossover, expressed in mask-fraction (noise-level) terms---not step index, which is a sampler hyperparameter that does not transfer across schedules---and reported as a \emph{distribution across samples}, since averaging can manufacture a smooth crossover that no individual trajectory exhibits.

\paragraph{Contributions.} This proposal contributes:
\begin{enumerate}
\item a mechanism-level hypothesis specific to \emph{discrete} diffusion---deterministic token commitment can absorb sub-margin token-identity perturbations, while confidence and scheduling remain separate error pathways (\S\ref{sec:hypothesis});
\item a measurement methodology that decomposes free-running divergence into per-step distortion and dynamical (non-)contractivity, avoiding the trajectory-forking confound that invalidates naive FP16-vs-quantized diffs (\S\ref{sec:exp1});
\item an experimental design in which every quantization claim carries a matched autoregressive control at the same scale, data, tokenizer, and recipe, so that only the \emph{gap-of-gaps} is interpreted (\S\ref{sec:exp2});
\item a versioned pre-registration plan for go/no-go criteria, including explicit kill conditions (\S\ref{sec:falsifiability}).
\end{enumerate}

\section{Background and Related Work}
\label{sec:related}

\paragraph{Quantized continuous diffusion.} For image diffusion, the question ``does quantization error accumulate over denoising steps?'' is answered: yes, non-uniformly by timestep, and the standard mitigations are timestep-aware calibration and mixed precision on sensitive steps (PTQ4DM~\cite{ptq4dm}, Q-Diffusion~\cite{qdiffusion}, TDQ~\cite{tdq}). This program does not re-derive these results; its contribution is confined to what differs in the discrete, masked-text setting.

\paragraph{Discrete diffusion language models.} MDLM~\cite{mdlm} provides small, well-documented masked-diffusion checkpoints with matched AR and SEDD baselines trained by the same authors on the same data. LLaDA~\cite{llada} and Dream demonstrate the paradigm at 7--8B scale.

\paragraph{Quantized dLLMs.} A fast-growing PTQ line now exists: systematic studies~\cite{dllmquant,qdlmstudy} document that AR-transfer PTQ degrades disproportionately on dLLMs and attribute this to timestep-dependent token distributions; Quant-dLLM~\cite{quantdllm} reaches 2-bit weights via masked-calibration simulation and mixed-precision allocation; STaR-Quant~\cite{starquant} names \emph{state-dependent activation disparity} and \emph{temporal error accumulation} explicitly and mitigates both inside calibration. Accumulation is therefore an \emph{observed and engineered-around} phenomenon. What this literature does not do---and what this program is built around---is (i) measure whether the underlying sampling dynamics contracts or expands injected error, separated from per-step distortion (\S\ref{sec:exp1}), (ii) run any comparison against a \emph{matched AR control} at the same precision, so none of it can attribute degradation to diffusion per se, (iii) study genuine-remasking samplers, and (iv) train a dLLM natively at low bit width: to our knowledge no dLLM QAT result exists at any precision.

\paragraph{Remasking.} A caution from code inspection motivates a key design choice: the well-known samplers of LLaDA and Dream are \emph{commit-and-freeze} despite ``remasking'' naming---LLaDA's low-confidence mode assigns $-\infty$ confidence to unmasked positions, so a committed token can never re-enter play. Commit-and-freeze removes revision, but it does not remove feedback: a wrong frozen token is still conditioning context for every later denoising call. Genuine remasking exists as an inference-time retrofit: ReMDM~\cite{remdm} re-opens committed tokens on any pretrained masked-diffusion checkpoint without retraining; trained self-correction (ProSeCo~\cite{proseco}) has since made revision a first-class training objective. Our primary sampler is therefore ReMDM; commit-and-freeze is measured on the same checkpoint as a secondary condition. This contrast tests the effect of revision policy with weights held fixed, but is not a feedback-only intervention: it can also change sampling distribution, compute, and quality. Those quantities are reported rather than attributed to feedback by assumption. Sampler implementations are validated against the sampler-correctness evaluations of Tang et al.~\cite{dllmsampler}.

\paragraph{Ternary models.} BitNet b1.58~\cite{bitnet} establishes both the promise (parity with FP16 from roughly 3B parameters) and the central methodological constraint: ternary PTQ fails on \emph{every} architecture, so native low-bit training is the path, and ternary-PTQ failure on a dLLM is uninformative unless it \emph{exceeds} the failure of a matched AR model quantized identically. Two adjacent results temper the priors: bottom-up QAT exploration~\cite{nielsen} finds ternary training on par with full precision even in small MLPs, GNNs, and encoder models---so small-scale ternary failure should not be assumed universal across architectures---and ternarized \emph{embedding spaces} have been reported beneficial in continuous-diffusion controllable text generation~\cite{qecdlm}, suggesting the embedding-exclusion default of \S\ref{sec:exp2} deserves its own ablation cell.

\section{Central Hypothesis: Discrete Commitment Can Suppress Some Error}
\label{sec:hypothesis}

The mechanism that distinguishes masked text diffusion from continuous diffusion is a \textbf{discrete token-commit decision}. With deterministic argmax decoding at a fixed position and fixed commit set, a logit perturbation smaller than the gap to the runner-up does not change token identity; a larger perturbation can flip a token discretely and, under commit-and-freeze, irreversibly. This is not total error absorption: the same perturbation can change confidence, position ranking, remasking, or stochastic sampling probabilities. The program tests whether identity preservation and revision outweigh those other pathways in end quality. It predicts:

\begin{enumerate}
\item \textbf{Conditional per-step tolerance.} At fixed state, fixed commit set, and deterministic decoding, text diffusion may preserve more token identities under ternary perturbation than a continuous output process preserves values.
\item \textbf{Per-mistake fragility.} A token flip can remain influential under commit-and-freeze, and confidence or schedule changes can fork a trajectory before any token flip is observed.
\item \textbf{Sampler dependence.} Revision policy changes the opportunities to correct a token flip and the sampling distribution. The commit/remask policy is therefore measured as a causal factor, not treated as a property of the denoiser alone.
\end{enumerate}

\section{Experimental Program}

The program is measurement-then-build: Experiments 0--1 establish whether and where ternary error dynamics are dangerous; Experiment 2 is the ternary training line the end-goal model grows out of.

\subsection{Experiment 0: PTQ Degradation with a Matched AR Control}
\label{sec:exp0}

\textbf{Objective:} establish the PTQ degradation profile of a small dLLM \emph{relative to a matched AR control} at each precision.

\textbf{Platform.} The primary model is MDLM-owt ($\sim$130M non-embedding parameters, GPT-2 tokenizer, OpenWebText)~\cite{mdlmcode}, chosen because (i) it is the checkpoint ReMDM was built against, and (ii) the same release ships an AR baseline trained by the same authors on the same data at the same scale in the same codebase---the matched control for free. A second scale point uses Tiny-A2D 0.5B/0.6B conversions~\cite{dllmframework}, whose AR \emph{parents} provide unusually tight controls. Larger open dLLMs (LLaDA-8B, Dream-7B, and an open-weight block-diffusion mixture-of-experts (MoE) model in the target model family) are measurement-only external-validity checks.

\textbf{Method.} Run FP16 baselines for both dLLM and AR control; apply naive PTQ to both at INT8, INT4, and ternary (precision pinned per result---INT4 findings do not transfer to ternary); apply a mask-aware PTQ variant to the dLLM (masked-calibration in the style of DLLMQuant and Quant-dLLM's MCS~\cite{quantdllm}); compare degradation \emph{deltas}: $(\text{dLLM}_{q} - \text{dLLM}_{fp16})$ versus $(\text{AR}_{q} - \text{AR}_{fp16})$.

\textbf{Interpretation.} If the dLLM's degradation materially exceeds the AR control's at matched precision, the refinement-dynamics worry is live. Comparable degradation supports ternary feasibility for dLLMs, but does not show that diffusion is \emph{more} tolerant. That positive claim requires a pre-specified gap-of-gaps favoring the dLLM, replicated at the second scale and without a quality or compute trade-off hidden in the sampler. Ternary PTQ failing on \emph{both} models is expected and does not answer the QAT question; the INT4$\to$ternary trend of the delta is only an early signal.

\subsection{Experiment 1: Error-Trajectory Dynamics}
\label{sec:exp1}

\textbf{Objective:} characterize how quantization error evolves under iterative refinement---the per-step distortion curve as a function of mask fraction, whether the dynamics contracts injected errors, and where absorption turns into amplification.

\textbf{The alignment confound.} The naive protocol---run FP16 and quantized denoisers from the same corrupted canvas and diff intermediate states---is confounded: as soon as the quantized model commits a different token at step $k$, the two trajectories are solving different problems at step $k{+}1$, and measured divergence conflates per-step distortion with trajectory forking. We therefore run three protocols and keep them separate:

\begin{description}
\item[Protocol A --- teacher-forced.] At every step, the quantized model receives the exact canvas the FP16 model produced, isolating per-step distortion at a fixed state. Two limits are stated up front: the protocol is off-policy (states are drawn from the FP16 trajectory distribution), and \emph{nothing can accumulate}---Protocol A cannot observe amplification by construction. Growth of Protocol A error with decreasing mask fraction is \emph{state-dependent sensitivity}, not feedback amplification; conflating the two invalidates the analysis.
\item[Protocol B --- free-running.] Both models run independently from the same start. This is the deployment condition, but its divergence mixes distortion with forking and cannot by itself prove expansion.
\item[Protocol C --- paired perturbation-injection stability.] Inject a controlled perturbation at mask fraction $\sigma$ (flip the $k$ least-confident committed tokens, or add calibrated pre-commit logit noise) and track token-Hamming distance to an unperturbed run. Run this paired experiment under both the FP16 and quantized transition maps, from both FP16 and quantized-trajectory states, with identical sampler randomness. Report recovery as a function of steps-since-injection and $\sigma$. This distinguishes state stability from weight-induced transition error; Protocol A calibrates the latter but does not make the two equivalent.
\end{description}

The measurements are complementary, not a proof by composition. Protocol A estimates conditional model distortion; Protocol C estimates state stability for the tested transition map; Protocol B measures deployed divergence. Their agreement supports a mechanistic account, while disagreement localizes a missing pathway. In particular, bounded A plus contractive C is evidence for robustness only after C is also measured under the quantized transition map and on quantized-state trajectories. If B diverges while A stays bounded and both FP16 and quantized C are contractive, the likely pathway is commit or schedule selection; only then is higher precision on the remasking/confidence path a justified mitigation.

\textbf{Metrics.} All trajectory metrics are functions of mask fraction and reported as per-sample distributions, never only means: committed-token agreement, top-two logit margins, confidence-ranking and remasking-selection disagreement, Hamming-decay curves after injection, and oscillation rate (remasking sampler only). Every divergence metric carries a locked evaluator and task score: generation is multimodal, so a diverged completion can be equally good---\textbf{a divergence finding with no quality drop is a fidelity result, not a failure}, and is labeled as such. Evaluator version, prompts, decoding randomness, sample count, and thresholds for ``grows,'' ``bounded,'' and ``contracts'' are committed in a timestamped configuration before runs.

\textbf{Expected findings.} (i) A crossover distribution in mask-fraction terms; if bimodal (early-fork vs.\ smooth-drift samples), the bimodality is the finding. (ii) Early-fork divergence implies a schedule/threshold fix; late smooth drift corroborated by Protocol C expansion implies genuine precision sensitivity near convergence. (iii) Mitigation priced honestly: higher precision below the crossover requires a resident higher-precision copy or on-the-fly dequantization, charged against the memory story---a crossover at 40\% mask fraction with a resident FP16 copy is not a ternary model.

\subsection{Experiment 2: Native Ternary Training and the Gap-of-Gaps}
\label{sec:exp2}

\textbf{Objective:} train the actual artifact---a natively ternary dLLM via quantization-aware training (QAT)---and test whether ternary hurts diffusion more, equally, or less than it hurts a matched AR model.

\textbf{The scale confound, stated up front.} BitNet's competitiveness threshold is roughly 3B parameters for decoder-only LLMs; at affordable scales ternary QAT is expected to underperform FP16 for reasons unrelated to diffusion, so a small ternary dLLM losing to its FP16 twin discriminates nothing. This prior is held loosely---ternary parity has been observed at small scale in non-decoder architectures~\cite{nielsen}---but either way the discriminating quantity is unchanged. The discriminating quantity is the \textbf{gap-of-gaps}:
\[
\underbrace{(\text{ternary dLLM} - \text{FP16 dLLM})}_{\text{diffusion ternary gap}} \quad \text{vs.} \quad \underbrace{(\text{ternary AR} - \text{FP16 AR})}_{\text{AR ternary gap}}
\]
at the same scale, tokenizer, data, token budget, and recipe. A comparable dLLM gap means diffusion adds no measured ternary fragility and the end-goal model is plausible wherever ternary AR is. A materially smaller dLLM gap is the positive result required to claim superior diffusion tolerance; a materially worse dLLM gap implicates refinement dynamics.

\textbf{Ablation grid} (each cell a separate QAT run, $\geq 3$ seeds; effects claimed only when between-cell differences exceed a pre-registered multiple of within-cell seed spread): all-ternary; all-ternary except embeddings/output head; all-ternary except mask embedding and remasking/confidence path; all-ternary except self-conditioning path; FP16 control. Per-module ablations are diagnostic only---quantization effects are non-additive, so all-ternary cells are never dropped.

\textbf{End-goal build.} The shipping artifact is an autoregressive-to-diffusion (A2D) conversion of a small open-weight AR model ($\sim$2--3B effective parameters, near BitNet's competitiveness threshold) with ternary QAT distilled from its FP16 parent, which doubles as both distillation teacher and matched AR control. The two distribution shifts are staged: convert to diffusion at FP16 first, verify against the AR parent, then ternarize with the FP16 diffusion model as teacher; the one-stage variant is an ablation, not the plan. The pipeline is debugged on 270M--1B pilots before the main run.

\textbf{Fragility points with defaults.} (i) The mask embedding and remasking/confidence path stay higher precision until ablations prove them safe---under commit-and-freeze this path is nearly the entire quantization question. (ii) Self-conditioning paths are tested separately. (iii) MoE routers stay FP16/INT8; routing errors are non-local.

\section{Systems and Memory Reality}
\label{sec:systems}

Two honesty constraints bound the deployment story. First, current GPUs have no native ternary matmul path; ternary is expected to be \emph{slower} than a good INT4 baseline until custom kernels prove otherwise, and serving begins with a standalone fixed-canvas inference loop, not an autoregressive-optimized stack like vLLM. Second, weight compression is the wrong lever if activations, embeddings, logits, and runtime buffers dominate: the popular headline ``a 25B ternary dLLM in 6\,GB'' is probably dead on arrival, and any claim about it must rest on a \emph{measured} peak-VRAM breakdown on the actual small model under test---weights (including any mixed-precision tail), embeddings/head, per-step activations, logits, and runtime buffers at the real canvas size, batch, and step count. That measurement is a required output of Experiments 0--1.

The useful deployment question is: at a fixed VRAM budget, does a ternary dLLM beat the best smaller INT4/FP16 autoregressive model on quality, latency, and reliability? At the scales this program can afford, dLLMs lose to AR models regardless of quantization, so the AR frontier is reported as \emph{context, not the kill gate}. If the answer stays negative at scale, the artifact ships as science, not as a serving strategy---an acceptable outcome, stated in advance.

\section{Falsifiability: Pre-Registration and Kill Criteria}
\label{sec:falsifiability}

Before data collection, all go/no-go thresholds are committed in versioned experiment configurations with a timestamp and repository revision; until then this section is a falsification plan, not a pre-registration. The program stops or re-scopes when any of the following holds: (i) Protocol C shows expansion over most of the schedule, Protocol A distortion is large there, \emph{and} the quality anchor drops past the declared threshold; (ii) quantizing the mask/remasking path changes commit decisions enough to miss quality thresholds; (iii) the ternary dLLM's gap to its FP16 twin exceeds the matched ternary AR gap by the pre-registered margin and the \S\ref{sec:exp2} mitigations do not close it; (iv) the crossover sits at high mask fraction, so the required higher-precision tail erases the memory advantage; (v) ternary inference remains slower than INT4/FP16 after reasonable kernel effort. Divergence-only signals with no quality drop never trigger kills, and losing to an \emph{unmatched} off-the-shelf AR model is context, never a kill. Notably, the null result at Experiment 0---dLLM and AR degrading comparably---does not kill the program; it re-scopes the science toward the argmax-absorption mechanism while the build proceeds on the strength of the ternary-AR analogy.

\section{Conclusion}

Diffusal should not be sold as ``ternary makes huge diffusion LLMs fit on tiny GPUs''---the arithmetic does not support that. The sharper, positive but falsifiable thesis is: \emph{can discrete commitment and revision make a dLLM at least as ternary-tolerant as a matched AR model, and can they make it more tolerant without moving cost or quality elsewhere?} A crossover in state stability is a possible explanation, not the expected answer. Token identity can be preserved below a margin, but confidence and schedule perturbations remain mechanisms to measure. If diffusion's ternary gap matches autoregression's, the result supports feasibility; only a replicated smaller gap supports the stronger absorption claim. If it is worse, the negative result is still real---iterative refinement is more precision-sensitive than its weight count suggests. Either way, the finding is trustworthy only if the transition maps are measured on their own trajectories, every quantization claim carries its matched AR control, and quality, memory, and compute are reported together.

\begin{thebibliography}{99}\small

\bibitem{bitnet} S. Ma et al. The Era of 1-bit LLMs: All Large Language Models are in 1.58 Bits. \emph{arXiv:2402.17764}, 2024.

\bibitem{mdlm} S. Sahoo et al. Simple and Effective Masked Diffusion Language Models. \emph{arXiv:2406.07524}, 2024.

\bibitem{llada} S. Nie et al. Large Language Diffusion Models. \emph{arXiv:2502.09992}, 2025.

\bibitem{dllmquant} C. Xu, D. Yang. DLLMQuant: Quantizing Diffusion-based Large Language Models. \emph{arXiv:2508.14090}, 2025.

\bibitem{qdlmstudy} H. Lin, H. Xu, Y. Wu, Z. Guo, R. Zhang, Z. Lu, Y. Wei, Q. Zhang, Z. Sun. Quantization Meets dLLMs: A Systematic Study of Post-training Quantization for Diffusion LLMs. \emph{arXiv:2508.14896}, 2025.

\bibitem{quantdllm} T. Zhang, Z. Li, X. Yan, H. Qin, Y. Guo, Y. Zhang. Quant-dLLM: Post-Training Extreme Low-Bit Quantization for Diffusion Large Language Models. \emph{arXiv:2510.03274}, ICLR 2026.

\bibitem{starquant} X. Yan, A. Wang, Z. Wan, X. Yu, I. Tsang. STaR-Quant: State-Time Consistent Post-Training Quantization for Diffusion Large Language Models. \emph{arXiv:2606.04945}, 2026.

\bibitem{nielsen} J. Nielsen, L. Galke, P. Schneider-Kamp. When Are 1.58 Bits Enough? A Bottom-up Exploration of Quantization-Aware Training with Ternary Weights. \emph{ICAART}, 2025. arXiv:2411.05882.

\bibitem{proseco} Y. Schiff, O. Belhasin, R. Uziel, G. Wang, M. Arriola, G. Turok, R. Zilberstein, M. Elad, V. Kuleshov. Learn from Your Mistakes: Self-Correcting Masked Diffusion Models. \emph{arXiv:2602.11590}, 2026.

\bibitem{qecdlm} C. Kang, X. Chen, Y. Hu, D. Novak. Quantized Embedding Vectors for Controllable Diffusion Language Models. \emph{TechRxiv preprint}, 2024.

\bibitem{ptq4dm} Y. Shang et al. Post-training Quantization on Diffusion Models. \emph{arXiv:2211.15736}, 2022.

\bibitem{qdiffusion} X. Li et al. Q-Diffusion: Quantizing Diffusion Models. \emph{arXiv:2302.04304}, 2023.

\bibitem{tdq} J. So et al. Temporal Dynamic Quantization for Diffusion Models. \emph{arXiv:2306.02316}, 2023.

\bibitem{remdm} G. Wang, Y. Schiff, S. Sahoo, V. Kuleshov. Remasking Discrete Diffusion Models with Inference-Time Scaling. \emph{arXiv:2503.00307}, 2025.

\bibitem{mdlmcode} MDLM checkpoints and matched AR/SEDD baselines. \url{https://github.com/kuleshov-group/mdlm}.

\bibitem{dllmframework} dLLM: Simple Diffusion Language Modeling (Tiny-A2D, AR$\to$diffusion conversion recipes). \emph{arXiv:2602.22661}, 2026. \url{https://github.com/ZHZisZZ/dllm}.

\bibitem{dllmsampler} L. Tang et al. Is Your Diffusion Sampler Actually Correct? A Sampler-Centric Evaluation of Discrete Diffusion Language Models. \url{https://github.com/LuhanTang/dllm_sampler}.

\end{thebibliography}

\end{document}
